% from main.tex
%%%%%%%%%%  scan idx 159  |  folio 146  %%%%%%%%%%
deux de gauche étant des homomorphismes de préfaisceaux d'anneaux), qui induisent, d'après (6.7), des isomorphismes sur les faisceaux associés. On va montrer que les faisceaux associés sont en fait canoniquement isomorphes au faisceau constant \(\mathbb Z_S\).

\textbf{Lemme 6.10.} \emph{Soient \((S,A)\) un topos localement annelé, \(U\) un objet de \(S\) distinct de l'objet vide, \(p,q\in\mathbb N\). Si l'on a \(A_U^p\simeq A_U^q\), alors \(p=q\).}

\emph{Preuve.} Par un argument similaire à celui donné dans la preuve de (2.15.1), on se ramène au cas où \((S,A)\) est le topos localement annelé défini par un préschéma, et dans ce cas l'assertion est triviale (prendre un point dans \(U\)).

\textbf{Définition 6.11.} On appelle rang, et l'on note \(\rg\) la \(S\)-fonction additive de \(\mathcal C_1\) dans \(\mathbb Z_S\) définie par
\[
  \rg(L)=p\quad\text{si } L\simeq A_X^p,
  \qquad X\ne\varnothing,
\]
\[
  \rg(L)=0\quad\text{si } L\in\Ob\mathcal C_{1X},
  \qquad X=\varnothing.
\]
En vertu de (6.7), la fonction rang se prolonge de manière unique en une \(S\)-fonction additive de \(\Parf(S)\) dans \(\mathbb Z_S\) (que nous appellerons encore rang et noterons \(\rg\)).

\textbf{Proposition 6.12.} \textbf{a)} Pour \(L,M\in\Ob\Dpx^-(X)_{\mathrm{parf}}\) (\(X\in\Ob S\)), on a \(\rg(L\overset{\mathrm L}{\otimes} M)=\rg(L)\,\rg(M)\).

\textbf{b)} Soit \(u:(S',A')\longrightarrow(S,A)\) un morphisme de topos localement annelés. Pour \(M\in\Ob\Dpx^-(S)_{\mathrm{parf}}\), on a \(\rg(\mathrm L u^*(M))=u^*(\rg(M))\). En d'autres termes, l'homomorphisme rang : \(k(\Dpx^-(S)_{\mathrm{parf}})\longrightarrow \mathbb Z_S\) est un homomorphisme de préfaisceaux d'anneaux, compatible avec les images inverses.



% Additional notation for Exposé I
\providecommand{\Cpx}{\mathrm C}
\providecommand{\Kpx}{\mathrm K}
\providecommand{\Dpx}{\mathrm D}
\providecommand{\Ho}{\mathrm H}
\providecommand{\RHom}{\mathrm R\!\operatorname{Hom}}
\renewcommand{\Ob}{\operatorname{ob}}
\providecommand{\Hom}{\operatorname{Hom}}
\providecommand{\RHom}{\mathrm R\!\operatorname{Hom}}
\providecommand{\Ext}{\operatorname{Ext}}
\providecommand{\Tor}{\operatorname{Tor}}
\providecommand{\Tr}{\operatorname{Tr}}
\providecommand{\Add}{\operatorname{Add}}
\providecommand{\Cart}{\operatorname{Cart}}
\providecommand{\Parf}{\operatorname{Parf}}
\providecommand{\coh}{\operatorname{coh}}
\providecommand{\amp}{\operatorname{amp}}
\providecommand{\cl}{\operatorname{cl}}
\providecommand{\rg}{\operatorname{rg}}
\providecommand{\Id}{\operatorname{Id}}
\providecommand{\Ker}{\operatorname{Ker}}
\renewcommand{\Im}{\operatorname{Im}}
\providecommand{\cH}{\mathcal H}
\providecommand{\Zs}{\mathbb Z_S}
\providecommand{\A}{\mathcal A}
\providecommand{\cC}{\mathcal C}
\providecommand{\cCzero}{\mathcal C_0}
\providecommand{\otL}{\mathop{\otimes}^{\mathrm L}}
\providecommand{\tamp}{\operatorname{tor.amp}}
\providecommand{\pamp}{\operatorname{parf.amp}}
\providecommand{\pdim}{\operatorname{parf.dim}}
\providecommand{\bon}{\mathrm b}
\providecommand{\parf}{\operatorname{parf}}
\providecommand{\parff}{\operatorname{parf}}
\providecommand{\ac}{\mathrm{ac}}


% Additional notation for Exposé I, completion pages 154--166
\providecommand{\End}{\operatorname{End}}
\providecommand{\tr}{\operatorname{tr}}

